% DRAFT (Claude, 2026-11-10, Anders): user test of the literature-review skill,
% second version after Anders' judgement (see ../../.claude/skills/literature-review/lessons.md).
% Claims checked in paper.txt; passages listed in CHANGES.md.

\subsubsection{Sources of Uncertainty and Their Dependence}
\label{sec:lit_unc_sources}

Stochastic modelling can give a better representation of reality in hydropower
scheduling, because inflow and prices are not known when the water is
allocated. \citet{belsnes2016applying} compare stochastic and deterministic
strategies for Norwegian cascades and find that the stochastic strategy earns
more in every case. The main reason is that it sees a risk of spill that the
deterministic plan does not. How much it earns depends on the plant. For a pumped-storage
plant, \citet{lohndorf2013} find that the gain over a rolling deterministic
plan is small when the power capacity is small relative to the reservoir, but
several times larger when the plant can move its reservoir content quickly. A
plant with large pumping and generating power therefore has more to gain from
a stochastic policy than a conventional plant with the same reservoir.

Inflow is an important element of this uncertainty, since it determines how
much water is available over the season. It can be treated as known, as short-term models often do. Many studies that
schedule reservoir hydropower over a year or longer model it as stochastic,
however \citep{abgottspon2013stochastic,helseth2016,helseth2017iet,bore2026assessing}.
Inflow follows from the weather, which is why system models typically draw
inflow, temperature and wind together from historical weather years
\citep{helseth2017}. For a pumped-storage plant with natural inflow the same
applies, with the difference that the pump adds a second, controllable source
of water to the reservoir.

Prices are the second uncertainty, and which prices are modelled as stochastic
shapes what a model can say. Many owner-oriented studies treat energy prices
and inflow jointly as stochastic
\citep{lohndorf2013,helseth2016,kleiven2023,bore2026assessing}.
\citet{helseth2016} and \citet{helseth2017iet} add uncertain reserve-capacity
prices, which lets the policy weigh reserve sales against energy sales under
the conditions it will actually meet. Short-term bidding models often keep
only the energy price uncertain \citep{faria2011day,aasgard2020value}. Some
studies take prices as known and solve each price year separately
\citep{ak2019quantifying,chazarra2017}. Such models show how the result
depends on price conditions, but not how decisions adapt as prices become
known.

Ideally every price the plant trades on would be stochastic, since that is
the most realistic description of what it faces. Each added variable, however,
multiplies the number of states and the data needed to estimate the
transitions between them. \citet{helseth2016} need nine joint energy- and
reserve-price states and have only 80 simulated scenarios to estimate them
from. The number of stochastic variables is therefore in practice set by the
data available, not by what would be most realistic.

When several inputs are uncertain, their dependence is a further modelling
choice. Prices and inflow are negatively correlated in the Norwegian system,
and \citet{kleiven2023} show that ignoring this changes water values and
reservoir operation, although the loss in expected revenue was small in their
case. Energy and reserve prices may also be correlated, and
\citet{perezdiaz2020medium} find little difference between a joint and an
independent description for a Spanish plant. Some methods require
independence between the price and inflow processes
\citep{helseth2016}, which then has to be justified from the data. The effect
of dependence thus varies between cases, and it has not been examined for a
plant where prices also drive pumping. The present study keeps price and
inflow in joint states for that reason.

Overall, many studies treat prices and inflow as the uncertain inputs and keep
their number small for tractability, which fits the data a single plant can
draw on. Whether their dependence must be modelled depends on the case. For
pumped storage the price enters twice: as the revenue from generating and as the cost of pumping. How prices
are described over time may therefore carry more weight than the number of
uncertain inputs. The next subsection turns to when price information becomes
known and at what resolution.
